\section{The \texorpdfstring{$(k, \ell)$}{(k, l)}-pebble game}
\label{sec:pebble-game}

The $(k, \ell)$-count matroid of \cref{sec:count-matroid} is defined
combinatorially: a subset of $E(K_V)$ is independent precisely when
the induced subgraph is $(k, \ell)$-sparse. \Cref{sec:count-matroid}
established that this collection of subsets satisfies the matroid
axioms via an augmentation / $I$-component argument, but the proof
is not constructive: given an arbitrary edge set, it does not
exhibit a procedure to decide whether the set is independent.

The \emph{$(k, \ell)$-pebble game} of
Lee--Streinu~\cite{leeStreinu2008} (generalising the original
$(2, 3)$ algorithm of Jacobs--Hendrickson~\cite{jacobsHendrickson1997})
is the missing algorithmic side. It processes a list of candidate
edges one at a time, maintaining a partial orientation of the
accepted edges together with a distribution of \emph{pebbles} on
the vertices, and decides at each step whether to insert the
candidate or reject it. We formalise the basic version of the
algorithm (Lee--Streinu \S 3) and prove its correctness theorem
(Lee--Streinu \S 4, Theorem 8) in the matroidal regime $\ell < 2k$
that matches \cref{sec:count-matroid}, in a certificate form: on
success, the algorithm returns the partial orientation as a
witness of sparsity; on failure, it returns a vertex subset $V'$
together with a proof that $|\edgesIn{V'}| > k|V'| - \ell$,
witnessing non-sparsity.

\emph{Out of scope.} The \emph{component pebble game} of
Lee--Streinu \S 5 maintains the $I$-components of the partial
orientation explicitly to reject dependent edges in $O(1)$ rather
than $O(n)$ time, achieving an overall $O(n^2)$ bound; the
underlying data structure (\emph{union pair-find}) was developed
by Lee--Streinu--Theran~\cite{leeStreinuTheran2005}. The Henneberg-
sequence application of Lee--Streinu \S 6 generalises the
Henneberg construction (\cref{sec:henneberg}) from Laman
graphs to all $(k, \ell)$-tight graphs. Neither is in scope for
this chapter; both are candidates for follow-up phases.

\emph{Multigraphs.} Lee--Streinu prove the basic pebble game's
correctness (\S~3--4, Lemma~10 through Lemma~15) for multi-graphs
with loops and parallel edges, splitting on the size hypothesis
$n, n' \ge 1$ in the lower range $\ell \in [0, k]$ and $n, n' \ge 2$
in the upper range $\ell \in (k, 2k)$. Following the project
convention (\cref{sec:sparsity}), this chapter works in
$\mathrm{SimpleGraph}\,V$ and in the matroidal regime $\ell < 2k$.
Two simplifications come for free under that specialisation: the
loop-free regime forces $\mathrm{span}(v) = 0$ at every vertex, so
Invariant~(1) of L-S Lemma~10 collapses from $\mathrm{peb}(v) +
\mathrm{span}(v) + \mathrm{out}(v) = k$ to $\mathrm{peb}(v) +
\mathrm{out}(v) = k$ (\cref{lem:pebble-game-invariants}, item~(1));
and the L-S size split $n' \ge 1$ / $n' \ge 2$ unifies as a single
premise $\ell \le k|V'|$ on Invariants~(3) and~(4), since
$|V'| \ge 1$ together with $\ell \le k$ gives $\ell \le k|V'|$ and
$|V'| \ge 2$ together with $k < \ell < 2k$ gives $\ell < 2k \le
k|V'|$. The substantive content of the four invariants and of
L-S Lemma~13's reach-pebble argument is otherwise unchanged: in
$\mathrm{SimpleGraph}\,V$ a candidate edge $\{u, v\}$ has $u \ne v$,
so the reachability subset $V' := \mathrm{Reach}(u) \cup
\mathrm{Reach}(v)$ already has $|V'| \ge 2$ and the unified
$\ell \le k|V'|$ premise is discharged directly from $\ell < 2k$ at
the call site (\cref{lem:pebble-game-independent-brings-pebble-graph}).
The per-edge correctness iff of L-S Lemma~14 is encoded structurally
in the return type of \cref{def:tryAddEdge}, which has one
constructor per outcome: accept, proved $(k, \ell)$-sparse by
\cref{lem:pebble-game-tryAddEdgeWith-isSparse}, and reject, proved
non-sparse by \cref{lem:workhorseWitness-certifies}. The full
input-graph classification of L-S Lemma~15 (under-/well-/over-constrained
according as the input is sparse-but-not-tight / tight / spanning)
specialises in this chapter to the matroidal-regime certificate
form \cref{thm:pebble-game-correct}; the well- and over-constrained
sub-classifications need the $I$-component machinery of the
component pebble game and are deferred along with that algorithm.
L-S Corollary~11 (block characterisation) is similarly part of the
component pebble game's apparatus. The algorithm state
(\cref{def:partial-orientation} below) is decoupled from the shape
of the input edge container, so retargeting to a future
\texttt{MultiGraph} type is a relaxation of the no-parallel-pair
invariant on the partial orientation rather than a structural
rewrite.

\subsection{State and moves}
\label{sec:pebble-game-state}

The pebble game maintains a partial orientation of a subset of
$E(K_V)$ together with a non-negative integer of \emph{free
pebbles} at each vertex. Both pieces of data are determined by
the orientation: pebbles are bookkeeping for the per-vertex
out-degree budget.

\begin{definition}[Partial orientation]
  \label{def:partial-orientation}
  \lean{CombinatorialRigidity.PebbleGame.PartialOrientation,
    CombinatorialRigidity.PebbleGame.PartialOrientation.underline}
  \leanok
  Let $V$ be a finite type with decidable equality. A
  \emph{partial orientation} on $V$ is a finite set
  $D \subseteq V \times V$ of ordered pairs such that:
  \begin{enumerate}
    \item (no loops) $(v, v) \notin D$ for all $v \in V$;
    \item (no antiparallel pairs) $(u, v) \in D$ implies
      $(v, u) \notin D$.
  \end{enumerate}
  An edge $\{u, v\} \in \mathrm{Sym}_2 V$ is \emph{oriented in $D$}
  if either $(u, v) \in D$ or $(v, u) \in D$. The
  \emph{underlying edge set} of $D$, written $\underline{D} \subseteq
  E(K_V)$, is the set of oriented edges.
\end{definition}

\begin{definition}[Out-degree, pebble count, span]
  \label{def:pebble-counts}
  \lean{CombinatorialRigidity.PebbleGame.PartialOrientation.out,
    CombinatorialRigidity.PebbleGame.PartialOrientation.peb,
    CombinatorialRigidity.PebbleGame.PartialOrientation.span,
    CombinatorialRigidity.PebbleGame.PartialOrientation.outOn,
    CombinatorialRigidity.PebbleGame.PartialOrientation.pebOn}
  \leanok
  \uses{def:partial-orientation}
  Fix a partial orientation $D$ and an integer $k \ge 0$. For a
  vertex $v \in V$, write
  \[
    \mathrm{out}_D(v) := |\{u : (v, u) \in D\}|,
    \qquad
    \mathrm{peb}_{D,k}(v) := k - \mathrm{out}_D(v).
  \]
  Extend additively to subsets $V' \subseteq V$:
  $\mathrm{peb}_{D,k}(V') := \sum_{v \in V'} \mathrm{peb}_{D, k}(v)$,
  and write
  \[
    \mathrm{span}_D(V') := |\{(u, w) \in D : u, w \in V'\}|,
    \qquad
    \mathrm{out}_D(V') := |\{(u, w) \in D : u \in V',\ w \notin V'\}|.
  \]
  The pebble count is non-negative provided $\mathrm{out}_D(v) \le k$
  for all $v$, which is preserved by the algorithm
  (\cref{def:tryAddEdge}) as a structural invariant.
\end{definition}

\begin{definition}[Path-reversal move]
  \label{def:path-reversal}
  \lean{CombinatorialRigidity.PebbleGame.PartialOrientation.reverse,
    CombinatorialRigidity.PebbleGame.PartialOrientation.underline_reverse_eq}
  \leanok
  \uses{def:partial-orientation}
  Let $D$ be a partial orientation and $u_0 \to u_1 \to \cdots \to
  u_m$ a directed path in $D$ (a sequence of vertices with
  $(u_i, u_{i+1}) \in D$ for $0 \le i < m$, all $u_i$ distinct).
  The \emph{reversal} of this path replaces $D$ by
  $D \setminus \{(u_i, u_{i+1}) : 0 \le i < m\} \cup
  \{(u_{i+1}, u_i) : 0 \le i < m\}$. The result is again a partial
  orientation: no loops are introduced (the path is simple), and the
  no-antiparallel-pairs invariant is preserved by induction on $m$.
  Path reversal preserves $\underline{D}$, $\mathrm{out}_D(v)$ for
  all $v \neq u_0, u_m$, and decreases $\mathrm{out}_D(u_0)$ by $1$
  while increasing $\mathrm{out}_D(u_m)$ by $1$; equivalently, the
  pebble count $\mathrm{peb}_{D, k}$ rises by $1$ at $u_0$ (under
  the precondition $\mathrm{out}_D(u_0) \le k$ from
  \cref{lem:pebble-game-invariants}~(1)) and drops by $1$ at $u_m$
  (under $\mathrm{out}_D(u_m) < k$, the precondition for the
  reversal move to be valid). The invariant-preservation half,
  span invariance ($\mathrm{span}_D(V')$ unchanged on every $V'$),
  the per-vertex \texttt{out}/\texttt{peb} effects, and the
  subset-level $\mathrm{peb}_{D, k}(V') + \mathrm{out}_D(V')$
  invariance under reversal are formalised. The last is routed
  through the algebraic identity $\mathrm{peb}_{D, k}(V') +
  \mathrm{span}_D(V') + \mathrm{out}_D(V') = k|V'|$
  (\cref{lem:pebble-game-invariants}~(2)) applied to both $D$ and
  $D'$: since $\mathrm{span}_D(V')$ is invariant, $\mathrm{peb}(V')
  + \mathrm{out}(V')$ must match. (The subset-level
  $\mathrm{out}_D(V')$ itself is *not* invariant: it shifts by
  $[u_m \in V'] - [u_0 \in V']$, cancelling the corresponding shift
  in $\mathrm{peb}_{D, k}(V')$.)
\end{definition}

\begin{definition}[Arc-insertion move]
  \label{def:arc-insertion}
  \lean{CombinatorialRigidity.PebbleGame.PartialOrientation.addArc,
    CombinatorialRigidity.PebbleGame.PartialOrientation.underline_addArc}
  \leanok
  \uses{def:partial-orientation}
  Let $D$ be a partial orientation and $(u, v)$ an ordered pair with
  $u \ne v$ and $(v, u) \notin D$. The \emph{arc-insertion} of
  $(u, v)$ replaces $D$ by $D \cup \{(u, v)\}$. The result is again
  a partial orientation: $u \ne v$ rules out a self-loop, and
  $(v, u) \notin D$ rules out an antiparallel pair (no other pair is
  changed). The per-vertex effect, assuming additionally $(u, v)
  \notin D$ so that the insertion is a genuine extension, is to
  raise $\mathrm{out}_D(u)$ by $1$ and leave $\mathrm{out}_D(x)$
  unchanged for all $x \ne u$; equivalently the pebble count
  $\mathrm{peb}_{D, k}(u)$ drops by $1$ (under
  $\mathrm{out}_D(u) < k$, the precondition for the insertion to be
  valid). The invariant-preservation half, per-vertex
  \texttt{out}/\texttt{peb} effects, and subset-level effects on
  $\mathrm{span}_D(V')$, $\mathrm{outOn}_D(V')$, and
  $\mathrm{peb}_{D, k}(V')$ are formalised. The subset-level shape
  is three independent unified additive identities keyed on the
  position of $u$, $v$ relative to $V'$: $\mathrm{span}_D(V')$ rises
  by $1$ iff $u, v \in V'$; $\mathrm{outOn}_D(V')$ rises by $1$ iff
  $u \in V'$ and $v \notin V'$; $\mathrm{peb}_{D, k}(V')$ drops by
  $1$ iff $u \in V'$. Combining them, $\mathrm{peb}_{D, k}(V') +
  \mathrm{out}_D(V')$ drops by $1$ iff both endpoints lie in $V'$
  and is otherwise unchanged --- the algebraic content of
  \cref{lem:pebble-game-invariants}~(3) at the insertion step.
\end{definition}

\subsection{Invariants}
\label{sec:pebble-game-invariants}

The four invariants of Lee--Streinu Lemma~10 are the soundness
backbone of the algorithm: they are maintained at every state
reachable from the initial empty orientation through edge
insertions and path reversals, and Invariant~(4) directly gives
the sparsity of $\underline{D}$.

\begin{definition}[Reachable orientation]
  \label{def:reachable-orientation}
  \lean{CombinatorialRigidity.PebbleGame.PartialOrientation.Reachable}
  \leanok
  \uses{def:partial-orientation, def:pebble-counts,
    def:path-reversal, def:arc-insertion}
  Fix integers $k, \ell \ge 0$. The set of $(k, \ell)$-\emph{reachable}
  partial orientations on $V$ is the smallest collection containing the
  empty orientation and closed under:
  \begin{itemize}
    \item \emph{Path reversal} along a simple directed path
      $u_0 \to \cdots \to u_m$ in $D$ with $\mathrm{out}_D(u_m) < k$
      (\cref{def:path-reversal}).
    \item \emph{Accepted arc insertion} of an ordered pair $(u, v)$
      with $u \ne v$, $(u, v) \notin D$, $(v, u) \notin D$,
      $\mathrm{out}_D(u) < k$, and the threshold condition
      $\mathrm{peb}_{D, k}(u) + \mathrm{peb}_{D, k}(v) \ge \ell + 1$
      (\cref{def:arc-insertion}).
  \end{itemize}
  Lee--Streinu Lemma 10 (\cref{lem:pebble-game-invariants} below)
  bundles the four invariants that hold across every reachable
  orientation.
\end{definition}

\begin{lemma}[Pebble-game invariants]
  \label{lem:pebble-game-invariants}
  \lean{CombinatorialRigidity.PebbleGame.PartialOrientation.Reachable.out_le,
    CombinatorialRigidity.PebbleGame.PartialOrientation.Reachable.peb_add_out_eq,
    CombinatorialRigidity.PebbleGame.PartialOrientation.Reachable.pebOn_add_span_add_outOn,
    CombinatorialRigidity.PebbleGame.PartialOrientation.Reachable.pebOn_add_outOn_ge,
    CombinatorialRigidity.PebbleGame.PartialOrientation.Reachable.span_add_le}
  \leanok
  \uses{def:partial-orientation, def:pebble-counts,
    def:reachable-orientation, def:arc-insertion, def:path-reversal}
  Let $V$ be a finite type, $k, \ell$ integers with $\ell < 2k$,
  and $D$ a $(k, \ell)$-reachable partial orientation on $V$
  (\cref{def:reachable-orientation}). Then for every $v \in V$ and
  every $V' \subseteq V$ with $\ell \le k|V'|$ (which is implied
  by $|V'| \ge 1$ when $\ell \le k$, and by $|V'| \ge 2$ when
  $k < \ell < 2k$):
  \begin{enumerate}
    \item $\mathrm{peb}_{D, k}(v) + \mathrm{out}_D(v) = k$.
    \item $\mathrm{peb}_{D, k}(V') + \mathrm{span}_D(V') +
      \mathrm{out}_D(V') = k|V'|$.
    \item $\mathrm{peb}_{D, k}(V') + \mathrm{out}_D(V') \ge \ell$.
    \item $\mathrm{span}_D(V') \le k|V'| - \ell$ (formalised in the
      additive form $\mathrm{span}_D(V') + \ell \le k|V'|$, per the
      project's no-$\mathbb{N}$-subtraction-in-propositions
      convention).
  \end{enumerate}
\end{lemma}
\begin{proof}
  \leanok
  \uses{def:partial-orientation, def:pebble-counts,
    def:reachable-orientation, def:arc-insertion, def:path-reversal}
  Invariant (1) reduces to $\mathrm{out}_D(v) \le k$ under truncated
  $\mathbb{N}$-subtraction $\mathrm{peb}_{D, k}(v) := k -
  \mathrm{out}_D(v)$; the latter is proved by induction on the move
  sequence using the per-vertex shifts
  (\cref{def:path-reversal}, \cref{def:arc-insertion}) and the
  $\mathrm{out}_D(w) < k$ / $\mathrm{out}_D(u) < k$ preconditions
  from \cref{def:reachable-orientation}. Invariant (2) is then the
  unconditional algebraic identity $\sum_{v \in V'}(\mathrm{peb} +
  \mathrm{out}) = k|V'|$ combined with the partition
  $\sum \mathrm{out}_D(v) = \mathrm{span}_D(V') +
  \mathrm{out}_D(V')$ (\texttt{sum\_out\_eq\_span\_add\_outOn} in the
  Lean source). Invariant (4) is the algebraic rearrangement of (2)
  given (3). Invariant (3) is the substantive claim, proved by
  induction on the move sequence: the empty orientation has
  $\mathrm{peb}_{\emptyset, k}(V') + \mathrm{out}_{\emptyset}(V') =
  k|V'| \ge \ell$; path reversals leave the sum
  $\mathrm{peb}(V') + \mathrm{out}(V')$ invariant
  (\texttt{pebOn\_add\_outOn\_reverse\_eq}); accepted insertions
  shift the sum by $-1$ exactly when both endpoints lie in $V'$
  (\texttt{pebOn\_add\_outOn\_addArc\_add}), and on that branch the
  threshold $\mathrm{peb}(u) + \mathrm{peb}(v) \ge \ell + 1$
  combined with the pairwise lower bound $\mathrm{peb}(u) +
  \mathrm{peb}(v) \le \mathrm{peb}(V')$ supplies the extra unit.
  Compare Lee--Streinu Lemma 10.
\end{proof}

\subsection{The basic algorithm}
\label{sec:pebble-game-algorithm}

The algorithm processes a list of candidate edges one at a time.
For each edge it attempts to redirect pebbles in $D$ until both
endpoints together hold at least $\ell + 1$ free pebbles, then
inserts and consumes one pebble. If the pebble target cannot be
reached, the edge is rejected and the algorithm halts with a
failure witness.

\begin{definition}[Try-reach-pebble: DFS with path reversal]
  \label{def:tryReachPebble}
  \lean{CombinatorialRigidity.PebbleGame.PartialOrientation.tryReachPebbleWith,
    CombinatorialRigidity.PebbleGame.PartialOrientation.tryReachPebble,
    CombinatorialRigidity.PebbleGame.PartialOrientation.TryReachPebbleResult,
    CombinatorialRigidity.PebbleGame.PartialOrientation.TryReachPebbleResult.newOrient}
  \leanok
  \uses{def:partial-orientation, def:pebble-counts, def:path-reversal,
    thm:reachable-finding-correct}
  Given a partial orientation $D$, a source vertex $v$, and a
  predicate $P$ on $V$, the procedure $\mathrm{tryReachPebble}\,
  D\,v\,P$ runs $\mathrm{reachableFinding}$
  (\cref{def:reachable-finding}) with relation $R$ taken to be the
  outgoing-arc relation of $D$. On success it receives a witness
  walk $p$ from $v$ to a vertex $w$ with $P(w)$
  (\cref{thm:reachable-finding-correct}) and reverses every arc of
  $p$ in $D$ via \cref{def:path-reversal}, returning the updated
  $(D', w)$. On failure it propagates $\mathrm{none}$. The
  reachability closure $\mathrm{Reach}_D(v)$ used in later lemmas
  coincides with the visited set of a fully-exhausted DFS run.

  The implementation is split into a computable workhorse
  \texttt{tryReachPebbleWith}, which takes a caller-supplied
  out-neighbour enumeration $\mathtt{succ} : V \to \mathrm{List}\,V$
  together with a proof that it agrees with $D$'s arc relation, and a
  convenience wrapper \texttt{tryReachPebble} plugging in $D$'s own
  out-neighbour list. (The wrapper is \texttt{noncomputable} only
  because mathlib's \texttt{Finset.toList} is; the computability gap
  is closed in \cref{sec:executable}.) The soundness and completeness
  theorems are stated on the workhorse and inherit to the wrapper.
\end{definition}

\begin{definition}[Try-add-edge]
  \label{def:tryAddEdge}
  \lean{CombinatorialRigidity.PebbleGame.PartialOrientation.tryAddEdgeWith,
    CombinatorialRigidity.PebbleGame.PartialOrientation.tryAddEdge}
  \leanok
  \uses{def:tryReachPebble, def:partial-orientation, def:arc-insertion}
  Given a partial orientation $D$, parameters $k, \ell$, and a
  candidate edge $\{u, v\} \notin \underline{D}$ (with $u \neq v$),
  the procedure $\mathrm{tryAddEdge}\,k\,\ell\,D\,\{u, v\}$
  repeatedly calls $\mathrm{tryReachPebble}\,D\,u$ and
  $\mathrm{tryReachPebble}\,D\,v$ with target predicate
  $w \mapsto \mathrm{peb}_{D,k}(w) > 0 \wedge w \neq u \wedge w \neq v$
  (free pebble at $w$, not at $\{u, v\}$ themselves --- otherwise the
  path-reversal step would either be trivial or move a pebble between
  the two endpoints without increasing $\mathrm{peb}_{D,k}(u) +
  \mathrm{peb}_{D,k}(v)$). Each successful DFS at $u$ (resp.\ $v$)
  fires a path reversal that raises $\mathrm{peb}_{D,k}(u)$ (resp.\
  $\mathrm{peb}_{D,k}(v)$) by $1$, and the procedure recurses on the
  updated orientation. The loop terminates as soon as either:
  \begin{itemize}
    \item (accept) $\mathrm{peb}_{D, k}(u) + \mathrm{peb}_{D, k}(v)
      \ge \ell + 1$ is reached; insert the directed arc, choosing
      $(u, v)$ when $\mathrm{peb}_{D,k}(u) > 0$ and $(v, u)$ otherwise
      (at least one endpoint has a free pebble since the sum is
      $\ge 1$), and return the updated orientation $D'$; or
    \item (reject) both DFS attempts simultaneously fail (no free
      pebble reachable from $u$ or $v$).
  \end{itemize}
  Termination of the outer loop: $(\ell + 1) - (\mathrm{peb}_{D,k}(u)
  + \mathrm{peb}_{D, k}(v))$ strictly decreases per successful
  reversal --- the predicate's $w \neq u, v$ clauses guarantee the
  reversed path's head is $u$ (resp.\ $v$) and its tail is neither,
  so \texttt{peb\_reverse\_head} raises the head's pebble count by
  $1$ while \texttt{peb\_reverse\_of\_not\_endpoint} leaves the other
  endpoint's untouched.

  As with $\mathrm{tryReachPebble}$, the procedure is split into a
  computable workhorse \texttt{tryAddEdgeWith} (taking a caller-supplied
  out-neighbour enumeration, which must be supplied for every
  intermediate orientation since path reversal mutates the arc set)
  and a convenience wrapper \texttt{tryAddEdge} plugging in $D$'s own
  out-neighbour list. The procedure also takes a proof that $D$ is
  $(k, \ell)$-reachable (\cref{def:reachable-orientation}), which
  supplies the per-vertex out-degree bound $\mathrm{out}_D(\cdot) \le k$
  used throughout.

  The result is one of two outcomes: \emph{accept}, returning the
  extended orientation $D'$; or \emph{reject}, returning a blocking
  witness $w$ (\cref{def:workhorseWitness}) whose subset is the
  reach-union $V' := \mathrm{Reach}_D(u) \cup \mathrm{Reach}_D(v)$
  (\cref{def:reachClosureComputable}). The witness records the
  algorithmic invariants $\mathrm{out}_D(V') = 0$ and
  $\mathrm{peb}_{D, k}(V') \le \ell$ holding at the rejection point;
  these certify that the ambient graph fails $(k, \ell)$-sparsity on
  $V'$ (\cref{lem:workhorseWitness-certifies}), which the user-facing
  wrapper \cref{def:runPebbleGame} turns into the rejection branch's
  certificate.
\end{definition}

\begin{definition}[Blocking witness]
  \label{def:blockingWitness}
  \uses{def:isSparse}
  Given a simple graph $G : \mathrm{SimpleGraph}\,V$ and parameters
  $k, \ell$, a \emph{blocking witness} for $G$ is a finite subset
  $V' \subseteq V$ together with a proof that
  $|\edgesIn[G]{V'}| > k|V'| - \ell$. Equivalently, $G$ is not
  $(k, \ell)$-sparse on $V'$.
\end{definition}

\begin{definition}[Run-pebble-game]
  \label{def:runPebbleGame}
  \lean{CombinatorialRigidity.PebbleGame.PartialOrientation.runPebbleGameWith,
    CombinatorialRigidity.PebbleGame.PartialOrientation.runPebbleGame}
  \leanok
  \uses{def:tryAddEdge, def:workhorseWitness, def:partial-orientation}
  Given a simple graph $G$ on a finite vertex set $V$, integer
  parameters $k, \ell$ with $\ell < 2k$, and the matroidal-regime
  hypothesis $h_{\mathrm{mat}} : \ell < 2k$, the procedure
  $\mathrm{runPebbleGame}\,k\,\ell\,G\,h_{\mathrm{mat}} :
  \mathrm{PebbleGameResult}\,G\,k\,\ell$ folds $\mathrm{tryAddEdge}$
  over an arbitrary enumeration of $E(G)$ starting from the empty
  partial orientation, then packages the workhorse-level $\mathrm{Sum}$
  output as a $\mathrm{PebbleGameResult}\,G\,k\,\ell$ verdict
  (\cref{def:pebbleGameResult}). On success it returns
  $\mathrm{.accept}\,(D, h_{\underline{}}, h_{\mathrm{reach}})$ with
  $\underline{D} = E(G)$ and $D$ satisfying the four invariants of
  \cref{lem:pebble-game-invariants}; on failure it returns
  $\mathrm{.reject}\,(V', h_{\mathrm{size}}, h_{\mathrm{lt}})$ with the
  blocking subset $V'$ and the $G$-shaped sparsity-violation pair
  ($V' := w.V'$ for the underlying workhorse-level failure witness
  $w$ (\cref{def:workhorseWitness}), with the pair extracted via
  \cref{lem:workhorseWitness-certifies}). Termination of the fold:
  the number of unprocessed edges strictly decreases per iteration.

  As before, the procedure factors through a computable workhorse
  \texttt{runPebbleGameWith}, which takes the orientation state, a
  reachability witness, a caller-supplied out-neighbour enumeration,
  and an explicit list of candidate ordered pairs; the convenience
  wrapper \texttt{runPebbleGame} plugs in an enumeration of $E(G)$
  starting from the empty orientation. The workhorse stays computable
  whenever its inputs do, and IO-driven callers invoke it directly.
  Because the workhorse does not see the ambient graph $G$, it returns
  the raw accept/reject pair of \cref{def:tryAddEdge}; the wrapper
  repackages this as the $\mathrm{PebbleGameResult}$ verdict
  (\cref{def:pebbleGameResult}), discharging the $G$-dependent proof
  fields.

  Each fold step first checks that the candidate is well-formed — not
  a loop, and not already oriented in either direction in $D$ — and
  skips it as a no-op otherwise; on a well-formed candidate it invokes
  $\mathrm{tryAddEdge}$ and recurses on the result. The reachability
  property of the orientation is maintained across the fold, so it
  need not be rechecked at each step.
\end{definition}

\subsection{Soundness}
\label{sec:pebble-game-soundness}

The success branch of $\mathrm{runPebbleGame}$ is automatically a
witness of sparsity, by Invariant~(4) applied to $\underline{D} =
E(G)$. This is the easy half of Lee--Streinu Theorem 8
(corresponding to L-S Corollary 12). The proof routes through two
infrastructure invariants on the fold: $\mathrm{Reachable}\,k\,\ell$-
preservation (already established in
\cref{lem:pebble-game-invariants}) and underlying-edge tracking
(\cref{lem:runPebbleGameWith-underline-subset} below).

\begin{lemma}[Underlying-edge tracking through the fold]
  \label{lem:runPebbleGameWith-underline-subset}
  \lean{CombinatorialRigidity.PebbleGame.PartialOrientation.tryAddEdgeWith_underline,
    CombinatorialRigidity.PebbleGame.PartialOrientation.runPebbleGameWith_underline_subset}
  \leanok
  \uses{def:tryAddEdge, def:runPebbleGame, def:partial-orientation,
    def:path-reversal, def:arc-insertion}
  Let $D$ be a partial orientation, $k, \ell$ integers, and
  $\mathit{edges} : \mathrm{List}\,(V \times V)$ a list of candidate
  ordered pairs. If the algorithm accepts the edge list (with appropriate
  witness parameters and a successor function), returning $D'$ (an accepting
  orientation), then
  \begin{enumerate}
    \item $\underline{D} \subseteq \underline{D'}$ (no edge is ever
      removed by the move sequence); and
    \item $\underline{D'} \subseteq \underline{D} \cup
      \{s(u, v) : (u, v) \in \mathit{edges}\}$ (the only new edges
      come from list entries that hit the accept branch).
  \end{enumerate}
  In particular, $\underline{D'} \subseteq \underline{D} \cup
  \mathrm{Sym}_2(\mathit{edges})$, so when $D = \mathtt{empty}$ and
  $\mathit{edges}$ enumerates $E(G)$ for a simple graph $G$,
  $\underline{D'} \subseteq E(G)$ at termination.

  Per-step accept-branch glue: $\mathtt{tryAddEdgeWith}\,k\,\ell\,D\,
  u\,v\ldots = \mathrm{.inr}\,D'$ implies $\underline{D'} = \underline{D}
  \cup \{s(u, v)\}$. Reversal moves preserve $\underline{D}$
  (\cref{def:path-reversal}), so any number of internal DFS-plus-
  reversal recursions in $\mathtt{tryAddEdgeWith}$ leave the
  underlying-edge set unchanged; the final $\mathtt{addArc}$ adds
  exactly one edge (\cref{def:arc-insertion}). The bidirectional
  $\subseteq$ statement is enough for soundness on the $\subseteq E(G)$
  half; the converse $E(G) \subseteq \underline{D'}$ (every edge is
  accepted) requires a no-skip-fires argument that is stated and
  proved separately when needed by \cref{thm:pebble-game-soundness}.
\end{lemma}
\begin{proof}
  \leanok
  By structural induction on $\mathit{edges}$. Empty list: $D = D'$,
  both inclusions are trivial. Cons case $(u, v) :: \mathit{es}$:
  if the runtime checks fail, the step is a no-op and the IH on
  $\mathit{es}$ gives both inclusions directly (the upper bound's
  RHS only grows when we widen the candidate list). Otherwise
  $\mathtt{tryAddEdgeWith}$ fires; on success it returns
  $D_{\mathrm{mid}}$ with $\underline{D_{\mathrm{mid}}} = \underline{D}
  \cup \{s(u, v)\}$ (by induction on the five cases of
  $\mathrm{tryAddEdge}$'s recursion, with the two
  accept branches closing on \cref{def:arc-insertion}'s underlying-
  edge equation and the two DFS-success branches composing the
  reversal-preservation lemma with the IH; the both-DFS-fail branch
  is contradictory by $\mathrm{.inr} \ne \mathrm{.inl}$ on the
  $\mathrm{Sum}$ output). The IH on $\mathit{es}$ at $D_{\mathrm{mid}}$
  then gives both inclusions for the full list.
\end{proof}

\begin{lemma}[No skip fires; underlying-edge identity on the wrapper]
  \label{lem:runPebbleGame-underline-eq}
  \lean{CombinatorialRigidity.PebbleGame.PartialOrientation.runPebbleGameWith_mem_underline,
    CombinatorialRigidity.PebbleGame.PartialOrientation.runPebbleGameWith_underline_eq,
    CombinatorialRigidity.PebbleGame.PartialOrientation.runPebbleGame_edges_discharges}
  \leanok
  \uses{def:partial-orientation, def:tryAddEdge, def:runPebbleGame,
    def:reachable-orientation, lem:runPebbleGameWith-underline-subset,
    lem:pebble-game-invariants}
  Let $D$ be a $\mathrm{Reachable}\,k\,\ell$-orientation and
  $\mathit{edges} : \mathrm{List}\,(V \times V)$ a list of ordered
  pairs satisfying: no entry is a loop ($p_1 \ne p_2$); every
  $s(p_1, p_2)$ for $p \in \mathit{edges}$ lies outside $\underline{D}$;
  and the entries are pairwise $\mathrm{Sym}_2$-distinct. If
  the algorithm accepts this list (with appropriate witness parameters and
  a successor function), returning an accepting orientation $D'$,
  then every $p \in \mathit{edges}$ satisfies
  $s(p_1, p_2) \in \underline{D'}$.

  In particular, the workhorse
  $\mathrm{runPebbleGameWith}\,\mathrm{empty}\,k\,\ell\,\ldots\,\mathit{edges}
  = \mathrm{.inr}\,D'$ against an edge list whose $\mathrm{Sym}_2$-image
  is $E(G)$ for some $G : \mathrm{SimpleGraph}\,V$ yields $\underline{D'}
  = E(G)$; the math-layer verdict-bearing wrapper
  $\mathrm{runPebbleGame}\,G\,k\,\ell$ (\cref{def:runPebbleGame})
  consumes this identity to populate $\mathrm{.accept}$'s
  $h_{\underline{}}$ proof field.
\end{lemma}
\begin{proof}
  \leanok
  Structural induction on $\mathit{edges}$. Empty: vacuous. Cons case
  $(u, v) :: \mathit{es}$: the three runtime checks at $(u, v)$ all
  pass — loop-freeness from the input; $(u, v), (v, u) \notin D.\mathrm{arcs}$
  from $s(u, v) \notin \underline{D}$ via the membership characterisation
  of $\underline{D}$ — so $\mathtt{tryAddEdgeWith}$
  fires (the per-vertex out-degree bound $\forall x, D.\mathrm{out}\,x
  \le k$ from \cref{lem:pebble-game-invariants}~(1) follows from the
  reachability hypothesis on $D$ and is not rechecked per step);
  on success
  $\underline{D_{\mathrm{mid}}} = \underline{D} \cup \{s(u, v)\}$
  (\cref{lem:runPebbleGameWith-underline-subset}'s per-step glue)
  and $\mathrm{Reachable}\,k\,\ell\,D_{\mathrm{mid}}$
  (\cref{lem:pebble-game-invariants}'s preservation lemma). The
  remaining $\mathit{es}$ inherit the three input hypotheses against
  $D_{\mathrm{mid}}$ — freshness combines the original freshness with
  the pairwise hypothesis (each $s(p_1, p_2)$ for $p \in \mathit{es}$
  is distinct from $s(u, v)$). The IH on $\mathit{es}$ at
  $D_{\mathrm{mid}}$ handles every entry; $s(u, v) \in \underline{D'}$
  follows from the monotonicity half of
  \cref{lem:runPebbleGameWith-underline-subset} on $\mathit{es}$.

  For the wrapper, the input list
  $E(G).\mathrm{toList}.\mathrm{map}\,\mathrm{Quot.out}$ satisfies the
  three hypotheses: no loops ($\neg e.\mathrm{IsDiag}$ for $e \in E(G)$
  from $G$'s loop-freeness); freshness against $\underline{\mathtt{empty}}
  = \emptyset$; pairwise $\mathrm{Sym}_2$-distinctness from
  $\mathrm{Finset.nodup\_toList}$ via the round-trip $s((Quot.out\,e)_1,
  (Quot.out\,e)_2) = e$. Combining the converse subset with the
  $\subseteq$ half of \cref{lem:runPebbleGameWith-underline-subset}
  gives $\underline{D'} = E(G)$.
\end{proof}

\begin{lemma}[Span equals $\edgesIn{V'}$-count under the underline identity]
  \label{lem:span-eq-ncard-edgesIn}
  \lean{CombinatorialRigidity.PebbleGame.PartialOrientation.span_eq_ncard_edgesIn}
  \leanok
  \uses{def:partial-orientation, def:pebble-counts,
    lem:runPebbleGame-underline-eq, def:edgesIn}
  Let $D$ be a partial orientation with $\underline{D} = E(G)$ for some
  $G : \mathrm{SimpleGraph}\,V$. For every $V' \subseteq V$,
  \[
    \mathrm{span}_D(V') = |\edgesIn[G]{V'}|.
  \]
\end{lemma}
\begin{proof}
  \leanok
  The map $(a, b) \mapsto s(a, b)$ from $D.\mathrm{spanArcs}\,V'$ to
  $\edgesIn[G]{V'}$ is a bijection. Injectivity is the no-antiparallel
  invariant of $D$: if $s(a, b) = s(a', b')$ with both arcs in $D$,
  then either $(a, b) = (a', b')$ outright or $(a, b) = (b', a')$,
  but the second case puts $(a, b)$ and $(b, a)$ both in $D.\mathrm{arcs}$,
  violating no-antiparallel. Surjectivity uses the membership
  characterisation of $\underline{D}$: an edge
  $e = s(a, b) \in \edgesIn[G]{V'}$ sits in $\underline{D} = E(G)$, so
  some orientation $(a, b)$ or $(b, a)$ is in $D.\mathrm{arcs}$; both
  endpoints lie in $V'$ by the $\edgesIn{}$ membership.
\end{proof}

\begin{theorem}[Soundness of the pebble game]
  \label{thm:pebble-game-soundness}
  \lean{CombinatorialRigidity.PebbleGame.PartialOrientation.runPebbleGameWith_sound}
  \leanok
  \uses{def:runPebbleGame, lem:pebble-game-invariants, def:isSparse,
    lem:runPebbleGame-underline-eq, lem:span-eq-ncard-edgesIn}
  Let $V$ be finite, $G : \mathrm{SimpleGraph}\,V$, and $k, \ell$
  integers with $\ell < 2k$. If the workhorse-level
  $\mathrm{runPebbleGameWith}\,\mathrm{empty}\,k\,\ell\,\ldots = \mathrm{.inr}\,D$
  against an edge list whose $\mathrm{Sym}_2$-image is $E(G)$, then
  $G$ is $(k, \ell)$-sparse. (The statement is at the workhorse level;
  the user-facing wrappers \cref{def:runPebbleGame} inherit it through
  the verdict they return.)
\end{theorem}
\begin{proof}
  \leanok
  \uses{lem:pebble-game-invariants, lem:runPebbleGame-underline-eq,
    lem:span-eq-ncard-edgesIn}
  By \cref{lem:pebble-game-invariants} Invariant~(4) applied to
  $D$ at termination, every subset $V' \subseteq V$ in the relevant
  size range satisfies $\mathrm{span}_D(V') \le k|V'| - \ell$.
  Since $\underline{D} = E(G)$ (\cref{lem:runPebbleGame-underline-eq}),
  \cref{lem:span-eq-ncard-edgesIn} rewrites this as
  $|\edgesIn[G]{V'}| \le k|V'| - \ell$, i.e.\ $G$ is $(k, \ell)$-sparse
  on $V'$. (The matroidal-regime assumption $\ell < 2k$ is not used
  formally on the Lean side; the size-range constraint folds into
  the definition of $(k, \ell)$-sparsity as $\ell \le k|V'|$.)
\end{proof}

\subsection{Completeness}
\label{sec:pebble-game-completeness}

The harder direction shows that the algorithm cannot reject any
edge of a $(k, \ell)$-sparse graph. The argument routes through
Lee--Streinu Lemma~13: if the candidate edge $\{u, v\}$ is
independent (in the matroid sense), then a free pebble can be
brought to $u$ or $v$ from within $\mathrm{Reach}_D(u) \cup
\mathrm{Reach}_D(v)$ without disturbing the other endpoint's count.
Iterating produces the $\ell + 1$ pebbles required to insert.

\begin{lemma}[Independent edges yield reachable pebbles, algebraic form]
  \label{lem:pebble-game-independent-brings-pebble}
  \lean{CombinatorialRigidity.PebbleGame.PartialOrientation.Reachable.independent_brings_pebble}
  \leanok
  \uses{def:partial-orientation, def:pebble-counts,
    def:reachable-orientation, lem:pebble-game-invariants}
  Let $D$ be a $(k, \ell)$-reachable partial orientation, $u \ne v$
  vertices, and $V' \subseteq V$ a finset containing both $u$ and $v$
  with $\mathrm{out}_D(V') = 0$ (i.e.\ $V'$ is closed under $D$'s
  outgoing arcs --- the algorithmic instance has $V' = \mathrm{Reach}
  _D(u) \cup \mathrm{Reach}_D(v)$). Suppose
  \begin{enumerate}
    \item $\mathrm{span}_D(V') + 1 + \ell \le k|V'|$ (the algebraic
      form of the post-insertion sparsity bound on $V'$: adding the
      candidate edge $\{u, v\}$ to $\underline{D}$ keeps the
      $\le k|V'| - \ell$ count); and
    \item $\mathrm{peb}_{D, k}(u) + \mathrm{peb}_{D, k}(v) < \ell + 1$
      (the below-threshold condition at which the outer algorithm
      attempts DFS).
  \end{enumerate}
  Then there is a vertex $w \in V' \setminus \{u, v\}$ with
  $\mathrm{peb}_{D, k}(w) \ge 1$.
\end{lemma}
\begin{proof}
  \leanok
  \uses{lem:pebble-game-invariants}
  Invariant~(2) of \cref{lem:pebble-game-invariants} gives
  $\mathrm{peb}_{D, k}(V') + \mathrm{span}_D(V') + \mathrm{out}_D(V')
  = k|V'|$. With $\mathrm{out}_D(V') = 0$ this is
  $\mathrm{peb}_{D, k}(V') = k|V'| - \mathrm{span}_D(V')$, and the
  post-insertion sparsity bound rearranges to $\mathrm{peb}_{D, k}(V')
  \ge \ell + 1$. Subtracting the below-threshold hypothesis
  $\mathrm{peb}_{D, k}(u) + \mathrm{peb}_{D, k}(v) < \ell + 1$ leaves
  $\sum_{w \in V' \setminus \{u, v\}} \mathrm{peb}_{D, k}(w) > 0$, so
  some $w \in V' \setminus \{u, v\}$ has $\mathrm{peb}_{D, k}(w)
  \ge 1$.

  Compare Lee--Streinu Lemma 13. The L-S statement takes
  $V' = \mathrm{Reach}_D(u) \cup \mathrm{Reach}_D(v)$ and uses
  $\mathrm{fromEdgeSet}\,(\underline{D} \cup \{\{u, v\}\})$ being
  $(k, \ell)$-sparse as the hypothesis: that converts to the present
  algebraic form via the closure-of-reachability identity
  $\mathrm{out}_D(\mathrm{Reach}_D(\cdot)) = 0$ and the
  $\mathrm{span}_D \leftrightarrow |\edgesIn{}|$ bridge of
  \cref{lem:span-eq-ncard-edgesIn} (adjusted by $+1$ for the added
  $\{u, v\}$ edge inside $V'$, since $u, v \in V'$ by reflexivity of
  $\mathrm{Reach}_D$). The L-S proof's appeal to an $I$-component /
  block argument unfolds in the algebraic form to a direct
  Invariant~(2) application, since the sparsity bound is the only
  fact about the post-insertion graph the proof actually consumes.
\end{proof}

\begin{lemma}[Independent edges yield reachable pebbles, SimpleGraph form]
  \label{lem:pebble-game-independent-brings-pebble-graph}
  \lean{CombinatorialRigidity.PebbleGame.PartialOrientation.Reachable.independent_brings_pebble_simpleGraph_form}
  \leanok
  \uses{def:partial-orientation, def:pebble-counts,
    def:reachable-orientation, def:isSparse,
    lem:pebble-game-independent-brings-pebble,
    lem:span-eq-ncard-edgesIn}
  Let $D$ be a $(k, \ell)$-reachable partial orientation, $u \ne v$
  with $\{u, v\} \notin \underline{D}$, and $G$ a finite simple graph
  with $E(G) = \underline{D} \cup \{\{u, v\}\}$ that is
  $(k, \ell)$-sparse. Suppose the matroidal regime $\ell < 2k$ holds
  and the algorithm's below-threshold condition fires:
  $\mathrm{peb}_{D, k}(u) + \mathrm{peb}_{D, k}(v) < \ell + 1$. Then
  there is a vertex $w \in \mathrm{Reach}_D(u) \cup \mathrm{Reach}_D(v)$
  distinct from $u$ and $v$ with $\mathrm{peb}_{D, k}(w) \ge 1$.
\end{lemma}
\begin{proof}
  \leanok
  \uses{lem:pebble-game-independent-brings-pebble,
    lem:span-eq-ncard-edgesIn}
  Instantiate $V' := \mathrm{Reach}_D(u) \cup \mathrm{Reach}_D(v)$ in
  \cref{lem:pebble-game-independent-brings-pebble}. The reflexivity of
  $\mathrm{Reach}_D$ puts $u, v \in V'$, the closure of $V'$ under
  $D$'s out-arcs gives $\mathrm{out}_D(V') = 0$, and $|V'| \ge 2$
  combined with $\ell < 2k$ yields the size hypothesis $\ell \le k|V'|$.
  The sparsity bound $\mathrm{span}_D(V') + 1 + \ell \le k|V'|$ follows
  from $G$-sparsity at $V'$ together with the
  inclusion $(\text{Sym2-image of } D\text{-arcs spanning } V')
  \cup \{\{u, v\}\} \subseteq E_G[V']$, whose cardinality equals
  $\mathrm{span}_D(V') + 1$ by injectivity of $\mathrm{Sym2.mk}$ on
  arcs (\cref{lem:span-eq-ncard-edgesIn}'s injOn ingredient) and
  freshness of $\{u, v\}$ w.r.t.\ $\underline{D}$.
\end{proof}

\begin{definition}[Workhorse-level failure witness]
  \label{def:workhorseWitness}
  \lean{CombinatorialRigidity.PebbleGame.WorkhorseWitness}
  \leanok
  \uses{def:partial-orientation, def:pebble-counts,
    def:reachable-orientation}
  Fix parameters $k$ and $\ell$ and a vertex type $V$. A
  \emph{workhorse-level failure witness} $w =
  (D, V', (u, v), \dots)$ packages the algorithm-side data extracted at
  $\mathtt{tryAddEdgeWith}$'s case-5 branch (both DFS searches fail) into
  a record of the following fields, all defined relative to the
  partial orientation $D$ at the failure point:
  \begin{itemize}
    \item $D$ — the partial orientation at the failure point;
    \item $V' \subseteq V$ — the blocking subset (the algorithmic
      instance has $V' = \mathrm{Reach}_D(u) \cup \mathrm{Reach}_D(v)$);
    \item the pending edge $(u, v)$ with $u \ne v$, $u, v \in V'$;
    \item $\mathrm{out}_D(V') = 0$ — $V'$ is closed under $D$'s
      outgoing arcs;
    \item $\mathrm{peb}_{D, k}(V') \le \ell$ — the subset-level pebble
      bound, strictly stronger than the per-vertex below-threshold
      $\mathrm{peb}_{D, k}(u) + \mathrm{peb}_{D, k}(v) \le \ell$:
      combined with the DFS-failure ``no free pebble outside $\{u, v\}$
      in $V'$'' assertion, every other $w \in V'$ has zero pebble count
      and the subset sum collapses;
    \item $\{u, v\} \notin \underline{D}$ — the candidate edge is fresh
      w.r.t.\ the running orientation's underline; and
    \item $D$ is $(k, \ell)$-reachable
      (\cref{def:reachable-orientation}).
  \end{itemize}
  The witness is deliberately $G$-free: the algorithm layer does not
  reference an ambient graph $G : \mathrm{SimpleGraph}\,V$, so the
  witness records only algorithm-shaped data. The $G$-shaped
  sparsity-violation certificate is recovered at the wrapper layer by
  \cref{lem:workhorseWitness-certifies} once the bridge hypotheses
  $\{u, v\} \in E(G)$ and $\underline{D} \subseteq E(G)$ are supplied.
\end{definition}

\begin{lemma}[Workhorse witness certifies non-sparsity]
  \label{lem:workhorseWitness-certifies}
  \lean{CombinatorialRigidity.PebbleGame.WorkhorseWitness.certifies_against}
  \leanok
  \uses{def:workhorseWitness, lem:pebble-game-invariants,
    lem:span-eq-ncard-edgesIn}
  Let $w$ be a workhorse-level failure witness
  (\cref{def:workhorseWitness}) at parameters $(k, \ell)$ on a vertex
  type $V$, in the matroidal regime $\ell < 2k$. For any finite simple
  graph $G : \mathrm{SimpleGraph}\,V$ such that
  $\{w.uv_1, w.uv_2\} \in E(G)$ and $\underline{w.D} \subseteq E(G)$,
  the witness's blocking subset $V' := w.V'$ certifies non-sparsity:
  \[
    \ell \le k |V'| \quad \text{and} \quad k |V'| < |\edgesIn{V'}| +
    \ell.
  \]
\end{lemma}
\begin{proof}
  \leanok
  \uses{lem:pebble-game-invariants, lem:span-eq-ncard-edgesIn}
  The size hypothesis $\ell \le k |V'|$ follows from $|V'| \ge 2$ (the
  two distinct endpoints $u, v$ both sit in $V'$) and the matroidal
  regime $\ell < 2k$. Invariant~(2) of
  \cref{lem:pebble-game-invariants} on $V'$ with $\mathrm{out}_D(V') =
  0$ gives $\mathrm{peb}_{D, k}(V') + \mathrm{span}_D(V') = k |V'|$,
  so the witness's bound $\mathrm{peb}_{D, k}(V') \le \ell$ rearranges
  to $k |V'| \le \mathrm{span}_D(V') + \ell$. The
  $+1$ bridge of \cref{lem:span-eq-ncard-edgesIn} (in the variant
  used by \cref{lem:pebble-game-independent-brings-pebble-graph}'s
  proof) under freshness of $\{w.uv_1, w.uv_2\}$ w.r.t.\
  $\underline{D}$ and the subset hypothesis
  $\underline{D} \subseteq E(G)$ gives $\mathrm{span}_D(V') + 1 \le
  |\edgesIn{V'}|$. Combining, $k |V'| + 1 \le |\edgesIn{V'}| + \ell$,
  i.e., $k |V'| < |\edgesIn{V'}| + \ell$.
\end{proof}

\begin{lemma}[Pebble game accept implies independent]
  \label{lem:pebble-game-tryAddEdgeWith-isSparse}
  \lean{CombinatorialRigidity.PebbleGame.PartialOrientation.tryAddEdgeWith_isSparse}
  \leanok
  \uses{def:tryAddEdge, def:reachable-orientation, def:isSparse,
    lem:pebble-game-invariants, lem:span-eq-ncard-edgesIn}
  Let $D$ be a $(k, \ell)$-reachable partial orientation, $u \ne v$
  a candidate edge, and $G$ a finite simple graph with
  $E(G) = \underline{D} \cup \{\{u, v\}\}$. If
  $\mathrm{tryAddEdge}\,k\,\ell\,D\,\{u, v\}$ returns $\mathrm{.inr}\,D'$
  (accept), then $G$ is
  $(k, \ell)$-sparse.
\end{lemma}
\begin{proof}
  \leanok
  \uses{lem:pebble-game-invariants, lem:span-eq-ncard-edgesIn}
  A successful insertion produces $D'$ with $\underline{D'} =
  \underline{D} \cup \{\{u, v\}\} = E(G)$ (the underline identity is
  preserved by path-reversal and grows by the inserted edge on accept).
  The recursive structure of \cref{def:tryAddEdge} respects
  \cref{def:reachable-orientation}: each path-reversal and each accept
  step matches a constructor, so $D'$ is reachable. Invariant~(4) of
  \cref{lem:pebble-game-invariants} on $D'$ gives
  $\mathrm{span}_{D'}(V') + \ell \le k|V'|$ for every $V'$ satisfying
  the size hypothesis; \cref{lem:span-eq-ncard-edgesIn} under the
  underline identity rewrites the left-hand side as
  $|\edgesIn[G]{V'}|$, yielding $(k, \ell)$-sparsity. Compare
  Lee--Streinu Lemma 14.
\end{proof}

\subsection{Workhorse-level correctness theorem}
\label{sec:pebble-game-correctness}

This is the certificate form of Lee--Streinu Theorem~8 in the
matroidal regime. It is stated at the workhorse layer: the pebble
game from the empty orientation against an edge list enumerating
$E(G)$ accepts (returns $\mathrm{.inr}\,D$) exactly when $G$ is
$(k, \ell)$-sparse, and on rejection (the $\mathrm{.inl}\,w$ branch)
carries the blocking-witness data inline, with the sparsity-violation
certificate for $G$ recovered by \cref{lem:workhorseWitness-certifies}.
The user-facing wrappers inherit this through the verdict they return;
the verdict-form restatement is \cref{thm:pebbleGameResult-isAccept-iff}.

\begin{theorem}[Pebble game correctness, workhorse certificate form]
  \label{thm:pebble-game-correct}
  \lean{CombinatorialRigidity.PebbleGame.PartialOrientation.runPebbleGameWith_correct}
  \leanok
  \uses{def:runPebbleGame, def:workhorseWitness, def:isSparse,
    thm:pebble-game-soundness, lem:runPebbleGame-underline-eq,
    lem:workhorseWitness-certifies}
  Let $V$ be finite, $G : \mathrm{SimpleGraph}\,V$, and $k, \ell$
  integers with $\ell < 2k$. Then $G$ is $(k, \ell)$-sparse if and
  only if there exists a partial orientation $D$ with
  $\mathrm{runPebbleGameWith}\,\mathrm{empty}\,k\,\ell\,\ldots =
  \mathrm{.inr}\,D$ against an edge list whose $\mathrm{Sym}_2$-image
  is $E(G)$. Concretely:
  \begin{enumerate}
    \item If the workhorse returns $\mathrm{.inr}\,D$, then $G$ is
      $(k, \ell)$-sparse and $\underline{D} = E(G)$.
    \item If the workhorse returns $\mathrm{.inl}\,w$ for some
      $w : \mathrm{WorkhorseWitness}\,k\,\ell\,V$, then $w.V'$
      certifies $\ell \le k|w.V'| \wedge k|w.V'| <
      |\edgesIn[G]{w.V'}| + \ell$ via
      \cref{lem:workhorseWitness-certifies} --- a blocking witness
      for $G$.
  \end{enumerate}
\end{theorem}
\begin{proof}
  \leanok
  \uses{thm:pebble-game-soundness, lem:runPebbleGame-underline-eq,
    lem:workhorseWitness-certifies}
  Part~(1) combines \cref{thm:pebble-game-soundness} (sparsity from
  Invariant~(4) on the final state) with
  \cref{lem:runPebbleGame-underline-eq} ($\underline{D} = E(G)$
  through the no-skip-fires fold).

  Part~(2) extracts the witness directly from the algorithm's
  structural output: the workhorse-level
  $\mathtt{runPebbleGameWith\_witness\_bridges}$ lemma threads the
  fold's history to establish $s(w.uv_1, w.uv_2) \in E(G)$ (the
  failure-pending edge sits in $G$, since $w.uv \in \mathrm{edges}$
  and $\mathrm{edges}$'s $\mathrm{Sym}_2$-image is $E(G)$) and
  $w.D.\underline{} \subseteq E(G)$ (the orientation-at-failure's
  underline lies inside $G$, via accept-step monotonicity from the
  empty input). Then \cref{lem:workhorseWitness-certifies} discharges
  the workhorse data into the $G$-shaped sparsity-violation pair.

  The biconditional comes from the dichotomy on the two outcomes:
  $\mathrm{.inr}$ closes via~(1); $\mathrm{.inl}$ closes via~(2) by
  contradicting the sparsity hypothesis at the witness $w.V'$. The
  failure witness is produced directly at the rejection point of
  \cref{def:tryAddEdge}, so no separate existence-form extraction is
  needed.
\end{proof}

\subsection{User-facing verdict}
\label{sec:pebble-game-verdict}

The workhorse-level correctness theorem
(\cref{thm:pebble-game-correct}) states an iff between sparsity and
the algorithm's raw accept/reject output. The user-facing wrappers
re-package that iff structurally: \cref{def:runPebbleGame} and
\cref{def:runPebbleGameExec} return a two-constructor inductive
$\mathrm{PebbleGameResult}\,G\,k\,\ell$ whose constructor \emph{is}
the certificate.

\begin{definition}[Pebble-game verdict]
  \label{def:pebbleGameResult}
  \lean{CombinatorialRigidity.PebbleGame.PebbleGameResult,
    CombinatorialRigidity.PebbleGame.PebbleGameResult.isAccept}
  \leanok
  \uses{def:partial-orientation, def:reachable-orientation, def:isSparse,
    def:edgesIn, def:workhorseWitness}
  Fix parameters $k$ and $\ell$ and a finite simple graph
  $G : \mathrm{SimpleGraph}\,V$. The \emph{pebble-game verdict}
  $\mathrm{PebbleGameResult}\,G\,k\,\ell$ is the two-constructor
  inductive
  \begin{itemize}
    \item $\mathrm{.accept}\,(D, h_{\underline{}}, h_{\mathrm{reach}})$:
      $D$ is a partial orientation with $\underline{D} = E(G)$ and
      $D$ is $(k, \ell)$-reachable
      (\cref{def:reachable-orientation});
    \item $\mathrm{.reject}\,(V', h_{\mathrm{size}}, h_{\mathrm{lt}})$:
      $V' \subseteq V$ is a blocking subset carrying the
      sparsity-violation pair
      \[
        \ell \le k|V'| \quad \text{and} \quad k|V'| <
        |\edgesIn[G]{V'}| + \ell.
      \]
  \end{itemize}
  The Boolean projection $\mathrm{.isAccept} :
  \mathrm{PebbleGameResult}\,G\,k\,\ell \to \mathbb{B}$ returns
  $\mathrm{true}$ on $\mathrm{.accept}$ and $\mathrm{false}$ on
  $\mathrm{.reject}$. The verdict's constructor \emph{is} the
  certificate: the iff $G\text{ is }(k,\ell)\text{-sparse} \iff
  \exists D,\, \mathrm{the\ algorithm\ accepts}$ is absorbed into the
  verdict's type. The two wrappers \cref{def:runPebbleGame} and
  \cref{def:runPebbleGameExec} build the verdict from the workhorse's
  raw output, populating the $\mathrm{.reject}$ proof fields via
  \cref{lem:workhorseWitness-certifies} (which consumes the
  matroidal-regime hypothesis $\ell < 2k$) and the $\mathrm{.accept}$
  fields via the soundness chain above.
\end{definition}

\begin{theorem}[Pebble game correctness, verdict form]
  \label{thm:pebbleGameResult-isAccept-iff}
  \lean{CombinatorialRigidity.PebbleGame.PartialOrientation.runPebbleGame_isAccept_iff,
    CombinatorialRigidity.PebbleGame.PartialOrientation.runPebbleGameExec_isAccept_iff}
  \leanok
  \uses{def:pebbleGameResult, thm:pebble-game-correct, def:isSparse}
  Let $V$ be finite, $G : \mathrm{SimpleGraph}\,V$, and $k, \ell$
  integers with $\ell < 2k$. Then $G$ is $(k, \ell)$-sparse if and
  only if the verdict-bearing wrapper
  $\mathrm{runPebbleGame}\,G\,k\,\ell\,h_{\mathrm{mat}}$ (resp.\ its
  exec-layer sibling $\mathrm{runPebbleGameExec}$) returns an
  $\mathrm{.accept}$ verdict — equivalently, its $\mathrm{.isAccept}$
  Boolean projection is $\mathrm{true}$.
\end{theorem}
\begin{proof}
  \leanok
  \uses{thm:pebble-game-correct}
  The verdict's $\mathrm{.isAccept}$ projection is $\mathrm{true}$
  exactly when the underlying workhorse-level
  $\mathrm{runPebbleGameWith}$ returned $\mathrm{.inr}\,D$; by
  \cref{thm:pebble-game-correct} this is equivalent to $(k, \ell)$-sparsity.
\end{proof}

\subsection{Matroidal-independence corollary}
\label{sec:pebble-game-matroid}

Translating the correctness theorem through the $(k, \ell)$-count
matroid identification of \cref{sec:count-matroid} gives the
algorithmic characterisation of independence in the count matroid:
the pebble game decides matroidal independence.

\begin{corollary}[Pebble-game certificate of matroidal independence]
  \label{cor:pebble-game-countMatroid-indep}
  \lean{SimpleGraph.countMatroid_indep_iff_runPebbleGame}
  \leanok
  \uses{thm:pebble-game-correct, thm:pebbleGameResult-isAccept-iff,
    def:countMatroid, thm:countMatroid-indep-iff, def:runPebbleGame}
  Let $V$ be finite and $k, \ell$ integers with $\ell < 2k$. For
  every $G : \mathrm{SimpleGraph}\,V$, the edge set $E(G)$ is
  independent in the $(k, \ell)$-count matroid
  $\mathrm{countMatroid}\,V\,k\,\ell$
  (\cref{def:countMatroid}) if and only if the verdict-bearing
  $\mathrm{runPebbleGame}\,G\,k\,\ell\,h_{\mathrm{mat}}$ returns
  $\mathrm{.accept}\,(D, \dots)$ — equivalently,
  $(\mathrm{runPebbleGame}\,G\,k\,\ell\,h_{\mathrm{mat}}).\mathrm{isAccept}
  = \mathrm{true}$.
\end{corollary}
\begin{proof}
  \leanok
  \uses{thm:pebble-game-correct, thm:countMatroid-indep-iff,
    thm:pebbleGameResult-isAccept-iff}
  By \cref{thm:countMatroid-indep-iff}, $E(G)$ is independent in
  the count matroid iff $G$ is $(k, \ell)$-sparse. By
  \cref{thm:pebbleGameResult-isAccept-iff}, the latter holds iff
  the verdict-bearing wrapper's $\mathrm{.isAccept}$ projection is
  $\mathrm{true}$.
\end{proof}

\emph{Special case: planar rigidity.} At $(k, \ell) = (2, 3)$,
\cref{cor:pebble-game-countMatroid-indep} composed with the
matroid identification of \cref{thm:rigidityMatroid-indep-iff-rowIndependent}
gives an algorithmic decision procedure for membership in the
planar rigidity matroid $\mathcal{R}(V)$
(\cref{def:rigidityMatroid}). Restricting further to spanning
subgraphs $H \subseteq G$ with $|E(H)| = 2|V| - 3$ decides the
Laman property (\cref{def:isLaman}). This is the Jacobs--Hendrickson
$(2, 3)$-pebble game~\cite{jacobsHendrickson1997} as the named
algorithmic specialisation of the present chapter's main theorem.
